\documentclass[11pt,a4paper]{article}

\usepackage[T1]{fontenc}

\usepackage{lmodern,amsmath,amssymb,microtype}

\usepackage[margin=26mm]{geometry}

\usepackage[hidelinks]{hyperref}

\usepackage{xurl}
\hypersetup{pdftitle={Math and Theory: collection overview},pdfauthor={Dev D. Goyal}}

\title{Math and Theory: collection overview}

\author{Dev D. Goyal}

\date{27 October 2026}

\begin{document}

\maketitle

\begin{abstract}

The catalogue contains nine result families and nine preprints in mathematics, theoretical computer science, and mathematical physics. Eight papers have complete Lean proofs within their stated formal scopes. The thermal proof covers finite spectral statements, without a formal stochastic long-time identification theorem. The Gyires preprint proves inequalities for structured matrices and develops related polynomial bounds. Its principal order-five claims, the all-order four-block theorem, and the exceptional-vector theorem have complete Lean proofs. The results comprise counterexamples, exact special cases, and sharp sequence bounds. The collection makes no claim of first priority. The papers have not been peer reviewed.

\end{abstract}

\tableofcontents

\section{How to use the collection}

The \href{https://github.com/dvgyl/math-and-theory/blob/main/CONTENTS.md}{manuscript catalogue} links each paper to its formal statements and independent computations. The \href{https://github.com/dvgyl/math-and-theory/blob/main/verification/README.md}{verification guide} records the commands and their results. Seven families share a Lake project pinned to Lean 4.35.0. Its 29 selected Comparator statements include equality, sharpness, and auxiliary positivity claims. The decision-tree and thermal packages use separate Lean 4.19.0 environments and a pinned Mathlib revision. The decision-tree package has 29 axiom reports of its own. This is a separate count from the shared Comparator selection.

\section{Mathematics}

\subsection{Real positive-definite counterexamples to extended q-permanent monotonicity}

For every real parameter $q>1$, some real symmetric positive-definite matrix has a negative $q$-permanent derivative at $q$. An explicit matrix of order five decreases from $q=14$ to $q=15$. The construction also gives correlation matrices with no zero entries and integer-entry examples. The dimension can depend on $q$. The real-symmetric question on $[-1,1]$ remains outside the result.

\par\smallskip

\noindent\textbf{Scope:} Counterexample and all-parameter extension.

\par\smallskip

\noindent\href{https://github.com/dvgyl/math-and-theory/blob/main/preprints/q-permanent/paper.pdf}{Manuscript}\quad\href{https://github.com/dvgyl/math-and-theory/blob/main/lean/docs/q-permanent.md}{Formal statement and prior work}.

\subsection{Gyires's permanent inequality for structured matrices}

Every nonnegative real circulant matrix $A$ of order five satisfies $24\operatorname{per}(AA^T)\leq625\operatorname{per}(A)^2$. Equality holds exactly for constant matrices. This proves the order-five circulant case of Gyires\textquotesingle s ratio inequality. A bound in sorted-gap coordinates has sharp coefficient $484$. The conjecture for arbitrary doubly stochastic matrices remains open.

\par\smallskip

\noindent\textbf{Scope:} Structured matrix inequalities and related polynomial bounds. The principal order-five claims, the all-order four-block theorem, and the exceptional-vector theorem have Lean proofs.

\par\smallskip

\noindent\href{https://github.com/dvgyl/math-and-theory/blob/main/preprints/gyires-circulant-five/paper.pdf}{Manuscript}\quad\href{https://github.com/dvgyl/math-and-theory/blob/main/lean/docs/gyires-circulant-five.md}{Formal statement and prior work}.

\subsubsection*{Broader arguments in the same preprint}

The manuscript also proves results for circulants of orders six and seven, with sharp sorted-gap margins. It treats $aI+bQ+cQ^2$ in every dimension, where $Q$ is any permutation matrix and the coefficients are nonnegative. Further results cover four-block doubly stochastic matrices with unequal block sizes and boundary entries, matrices with one exceptional column or row, and every rank-at-most-two doubly stochastic matrix through order fifteen.

The supporting sections prove structural, scalar, deletion, balancing, local polynomial, and configuration-model results. Exact counterexamples identify the limits of auxiliary assertions. The \href{https://github.com/dvgyl/math-and-theory/blob/main/research-notes/gyires-complete-support/README.md}{complete support record} gives the source map and exact reproduction commands.

The all-order four-block and exceptional-vector theorems have complete Lean proofs. Their inequalities and equality cases are among the 29 selected formal statements. The principal higher-order circulant, permutation-polynomial, and finite rank-two theorems remain incomplete in Lean. The all-order rank-two problem and the unrestricted Gyires conjecture remain unproved here.

\subsection{An entropy increase for a structurally disassortative DomiRank graph}

A connected graph with $14$ vertices and $32$ edges has negative degree--log-degree edge covariance. All resolvent scores are positive on $[251/1000,63/250]$ inside the stable domain. Its entropy nevertheless increases by more than $1/10000$ in natural-logarithm units. This refutes the source\textquotesingle s assertion for every positive-score interval.

\par\smallskip

\noindent\textbf{Scope:} Counterexample to an interval monotonicity conjecture.

\par\smallskip

\noindent\href{https://github.com/dvgyl/math-and-theory/blob/main/preprints/domirank-entropy/paper.pdf}{Manuscript}\quad\href{https://github.com/dvgyl/math-and-theory/blob/main/lean/docs/domirank-entropy.md}{Formal statement and prior work}.

\subsection{Diagonal concavity does not imply concavity for sums of elementary symmetric polynomials}

For $M>0$, let $f_M=1+M\sigma_1+\sigma_2+M\sigma_3+\sigma_4$ in four variables, where $\sigma_j$ is the elementary symmetric polynomial of degree $j$. The fourth root $f_M^{1/4}$ is concave on the positive diagonal exactly when $M\geq1$. It is concave on the full positive orthant exactly when $M=1$. Thus every $M>1$ disproves Lachaume\textquotesingle s criterion with strictly positive coefficients.

\par\smallskip

\noindent\textbf{Scope:} Counterexample and complete parameter classification.

\par\smallskip

\noindent\href{https://github.com/dvgyl/math-and-theory/blob/main/preprints/diagonal-concavity/paper.pdf}{Manuscript}\quad\href{https://github.com/dvgyl/math-and-theory/blob/main/lean/docs/diagonal-concavity.md}{Formal statement and prior work}.

\subsection{Sharp root log-concavity of alternating Motzkin sums}

Let $T_n=\sum_{j=0}^n(-1)^{n-j}M_j$, where $M_j$ is the Motzkin number. The sequence $T_n^{1/n}$ is strictly log-concave on its tail starting at five. The inequality reverses at center five, so no earlier tail is log-concave. This strengthens Zhao\textquotesingle s proposed tail starting at nine. The source recurrence and finite sum are linked in Lean.

\par\smallskip

\noindent\textbf{Scope:} Conjecture resolution with an optimal threshold.

\par\smallskip

\noindent\href{https://github.com/dvgyl/math-and-theory/blob/main/preprints/alternating-motzkin/paper.pdf}{Manuscript}\quad\href{https://github.com/dvgyl/math-and-theory/blob/main/lean/docs/alternating-motzkin.md}{Formal statement and prior work}.

\subsection{Counterexamples to a late-iterate ordering conjecture for generalized Hofstadter functions}

For every $k\geq4$, some input $n$ satisfies $F_k^{\circ(2k+1)}(n)>F_{k+1}^{\circ(2k+2)}(n)$. Here the superscript denotes repeated composition. The proof constructs inverse iterates with exact gap $k-3$. The witness $F_4^{\circ9}(132)=8>7=F_5^{\circ10}(132)$ gives the first instance. The result concerns infinitely many parameters. It does not assert infinitely many violating inputs for each fixed parameter.

\par\smallskip

\noindent\textbf{Scope:} Counterexamples for every parameter k$\geq$4.

\par\smallskip

\noindent\href{https://github.com/dvgyl/math-and-theory/blob/main/preprints/hofstadter-iteration/paper.pdf}{Manuscript}\quad\href{https://github.com/dvgyl/math-and-theory/blob/main/lean/docs/hofstadter-iteration.md}{Formal statement and prior work}.

\section{Theoretical computer science}

\subsection{Counterexamples to a best-of-\texorpdfstring{$n$}{n} win-rate bound at every budget}

For every $n\geq2$ and $\delta>0$, a finite positive reference distribution with largest atom at most $\delta$ violates the proposed best-of-$n$ divergence--win-rate curve. A four-response best-of-two example gives exact probabilities and a strict logarithmic certificate. The all-budget proof uses finite interval bounds and gives a quantitative positive gap. It does not determine the sharp replacement curve.

\par\smallskip

\noindent\textbf{Scope:} Disproof for every budget n$\geq$2 and every positive atom cap.

\par\smallskip

\noindent\href{https://github.com/dvgyl/math-and-theory/blob/main/preprints/best-of-n/paper.pdf}{Manuscript}\quad\href{https://github.com/dvgyl/math-and-theory/blob/main/lean/docs/best-of-n.md}{Formal statement and prior work}.

\subsection[Decision-tree covariance: logarithmic bounds and variance counterexamples]{Logarithmic influence bounds and variance counterexamples for Boolean decision-tree covariance}

Let $T$ be a finite Boolean decision tree under uniform independent bits. No coordinate repeats on a root-to-leaf path. Let $k_i(T)$ count all query nodes labelled $i$, and let $I_i(f)$ be the probability that flipping coordinate $i$ changes the output. The weighted sum of absolute child covariances satisfies
\[
\sum_v 2^{-d(v)}\left|\operatorname{Cov}(f_{v,+},f_{v,-})\right|
\leq \sum_i I_i(f)\log_2\max\{1,k_i(T)\}.
\]
Here $d(v)$ is the node depth, and the two child functions use the same input cube. The read number counts nodes in the entire tree, not queries along one path.

For a read-$k$ tree, the bound is $I(f)\log_2 k$. The dummy-prefix construction of Wan, Wright, and Wu attains equality when $k$ is a power of two. The logarithmic potential is due to Shalev. A four-variable counterexample has covariance $5/4$ and variance $3/4$. A nine-variable counterexample has covariance $4863/4096$ and variance $3007/4096$, with every query relevant. Both refute the coefficient-one logarithmic variance proposal for a given tree.

\par\smallskip
\noindent\textbf{Scope:} Covariance inequalities, an equality family, and exact finite witnesses. No Fourier-entropy theorem is formalized.

\par\smallskip
\noindent\href{https://github.com/dvgyl/math-and-theory/blob/main/verification/additional-proofs/readk-covariance/docs/manuscript.pdf}{Manuscript}\quad\href{https://github.com/dvgyl/math-and-theory/blob/main/verification/additional-proofs/readk-covariance/docs/related-work.md}{Prior work}\quad\href{https://github.com/dvgyl/math-and-theory/blob/main/verification/additional-proofs/readk-covariance/README.md}{Formal proof and reproduction}.

\section{Mathematical physics}

\subsection{A three-state thermal counterexample to a heat-current kurtosis conjecture}

A three-state rate model with two positive-temperature baths satisfies local detailed balance. Every total off-diagonal rate is positive. The tilted generator records heat transferred to the cold bath. Its dominant eigenvalue branch is four times continuously differentiable near zero and has a larger real part than every distinct complex eigenvalue. The exact derivative ratio is
\[
\frac{C_4}{C_2}=-\frac{3470947982622836}{176627429595846691}<0,
\]
where $C_n$ is the $n$th derivative at zero. This refutes Eq.~(19) of Ptaszy\'nski, \emph{Physical Review E} \textbf{106}, 024119 (2022), in the paper's finite thermal spectral formulation. The third derivative is positive, so the witness does not refute the separate skewness conjecture.

\par\smallskip
\noindent\textbf{Scope:} Finite rate matrices, local detailed balance, complex eigenvalues, and derivatives of a local dominant branch. The stochastic long-time theorem that identifies the branch with a cumulant generating function is not formalized. No minimal-state or empirical validation claim is made.

\par\smallskip
\noindent\href{https://github.com/dvgyl/math-and-theory/blob/main/verification/additional-proofs/thermal-kurtosis-counterexample/report.pdf}{Manuscript}\quad\href{https://github.com/dvgyl/math-and-theory/blob/main/verification/additional-proofs/thermal-kurtosis-counterexample/RELATED_WORK.md}{Prior work}\quad\href{https://github.com/dvgyl/math-and-theory/blob/main/verification/additional-proofs/thermal-kurtosis-counterexample/README.md}{Formal proof and reproduction}.

\section{Verification and authorship}

Lean checks the proofs against their formal statements. In the shared Lean 4.35.0 library, separate Challenge modules state the intended claims. Comparator checks these against their proved Solutions. The verification record identifies the completed runs of Lean, NanoDa, and con-ron for those seven families. The two standalone Lean 4.19.0 packages have their own source compilation, axiom, and arithmetic records. The shared Comparator results do not cover those packages. Independent exact and numerical programs test finite certificates and examples. These finite computations do not replace the general proofs.

\par

Dev D. Goyal directed the research and is responsible for publication.
AI tools contributed mathematical arguments, Lean formalizations, computations, and editing.
The \href{https://github.com/dvgyl/math-and-theory/blob/main/AUTHORSHIP.md}{authorship disclosure} identifies these tools and their roles.
These preprints have not been peer reviewed.

\end{document}
